\newcolumntype{P}[1]{>{\raggedright\arraybackslash}p{#1}}

{\small
\renewcommand{\arraystretch}{1.3}
\begin{longtable}{|P{2.0cm}|P{2.3cm}|P{2.5cm}|P{1.8cm}|P{4.3cm}|}
\caption{Summary of major related works in causal discovery, regime detection,
and causal clustering.}
\label{tab:related_works} \\
\hline
\textbf{Study} & \textbf{Task} & \textbf{Method} & \textbf{Setting} & \textbf{Main Finding / Gap} \\
\hline
\endfirsthead

\multicolumn{5}{c}%
{{\bfseries \tablename\ \thetable{} -- continued from previous page}} \\
\hline
\textbf{Study} & \textbf{Task} & \textbf{Method} & \textbf{Setting} & \textbf{Main Finding / Gap} \\
\hline
\endhead

\hline \multicolumn{5}{|r|}{{Continued on next page}} \\ \hline
\endfoot

\hline
\endlastfoot

Spirtes et al.\ \cite{spirtes2000causation}
& DAG learning from observational data
& PC and FCI constraint-based algorithms
& i.i.d., single environment
& Foundational framework for causal discovery. Assumes a single fixed causal model throughout. \\
\hline
Zheng et al.\ \cite{zheng2018dags}
& Continuous DAG structure learning
& NOTEARS: smooth algebraic acyclicity constraint
& i.i.d., single environment
& Enables gradient-based DAG learning at scale. The acyclicity constraint underpins DYNOTEARS and CASTOR. \\
\hline
Pamfil et al.\ \cite{pamfil2020dynotears}
& Time-series causal graph recovery
& DYNOTEARS: NOTEARS extended to contemporaneous and lagged edges
& Stationary time series
& Efficiently recovers window causal graphs. Assumes stationarity holds throughout the entire series. \\
\hline
Huang et al.\ \cite{huang2020causal}
& Causal discovery under heterogeneity
& CD-NOD: nonparametric method exploiting distribution shifts
& Multiple domains or time-indexed data
& Heterogeneity aids causal identification. Requires an auxiliary domain index or time stamp. \\
\hline
Peters et al.\ \cite{peters2016causal}
& Causal parent identification
& ICP: invariant prediction across environments
& Multiple labelled environments
& Exploits causal invariance effectively. Requires environment labels as a strict prerequisite. \\
\hline
Arjovsky et al.\ \cite{arjovsky2019irm}
& Out-of-distribution generalisation
& IRM: invariant risk minimisation
& Multiple labelled environments
& Encourages causally stable features. Guarantees weaken when the environment count is small. \\
\hline
Perry et al.\ \cite{perry2022causal}
& Multi-environment causal graph recovery
& Sparse Mechanism Shift hypothesis with MSS score
& Multiple labelled environments
& Full graph identifiable under sparse shifts. Environment membership must still be known. \\
\hline
Huang et al.\ \cite{huang2019specific}
& Causal clustering with shared and specific relations
& SSCM: mixture of causal mechanisms
& i.i.d., multiple regimes
& Distinguishes shared and regime-specific causal relations. Requires a pre-specified regime count. \\
\hline
Chen et al.\ \cite{chen2023ccsl}
& Causal structure learning from unknown environments
& CCSL: Bayesian nonparametric Causal CRP with variational inference
& i.i.d., unknown environments
& Does not require a regime count in advance. Identifiable under a linear non-Gaussian assumption. Primary baseline for this thesis. \\
\hline
Varambally et al.\ \cite{varambally2024mcd}
& Mixture of causal models from time series
& MCD: variational inference over mixtures of SCMs
& Temporal, autoregressive structure
& Jointly infers causal models and mixing probabilities. Relies on temporal ordering. \\
\hline
Rahmani and Frossard \cite{rahmani2025castor}
& Temporal causal regime discovery
& CASTOR: EM over piecewise-constant regimes with NOTEARS acyclicity
& Sequential time series
& Jointly learns regime boundaries and per-regime DAGs without a prior regime count. Relies on temporal contiguity unavailable in i.i.d.\ settings. \\
\hline
Cheng et al.\ \cite{cheng2025dycast}
& Continuously evolving causal structure
& DyCAST: Neural ODE constrained to DAG manifold
& Sequential time series
& Captures smooth causal transitions and scales beyond 50 nodes. Requires temporal ordering. \\
\hline
Li et al.\ \cite{li2025hcl}
& Unsupervised causal clustering from mixed-type data
& HCL: bi-directional iterative causal clustering (preprint)
& i.i.d., no environment labels
& No temporal ordering or regime count required. Scalability on large graphs has not yet been assessed. \\
\hline
Sadeghi et al.\ \cite{sadeghi2025causal}
& Causal discovery from nonstationary time series
& CD-NOTS: nonparametric extension of CD-NOD correcting lagged conditioning
& Non-stationary time series
& Uncovers nonlinear nonstationary causal relations missed by stationarity-assuming methods. \\
\hline
Mooij et al.\ \cite{mooij2020jci}
& Joint causal inference from multiple contexts
& JCI: context encoded as additional graph nodes
& Multiple pre-specified contexts
& Theoretically general and subsumes many special cases. Contexts must be labelled in advance. \\
\hline
Faltenbacher et al.\ \cite{faltenbacher2025internal}
& Assumption violation detection
& Internal incoherency scores for constraint-based algorithms
& i.i.d., single environment
& Flags assumption violations without ground truth. Not yet adapted to causal clustering settings. \\
\hline

\end{longtable}
}
